\emph{(i)} For the derived metric $d\tau^{2} - \sum_{j}\tau^{2p_{j}}(dx^{j})^{2}$ the volume element is $\sqrt{|g|} = \tau^{\sum_{j}p_{j}} = \tau$. The massless minimally coupled field equation $|g|^{-1/2}\partial_{\mu}(|g|^{1/2}g^{\mu\nu}\partial_{\nu}\phi) = 0$ for a mode $e^{I_{4}kx^{i}}$ (numbers are read in $\mathrm{span}\{1, I_{4}\}$), constant along the other directions (for a hidden direction this is its zero mode), reads $\tau^{-1}\partial_{\tau}(\tau\dot\phi) + k^{2}\tau^{-2p_{i}}\phi = 0$, the stated equation. On a vacuum Kasner point the scalar contraction of the field strength vanishes (the derived metric is Ricci flat, Lemma \claimref{interior.kasner-form-interior} and Fact \claimref{sheaf.kasner-data-interior} for the Schwarzschild interior, and every point of the circle), so a coupling $\xi R\phi^{2}$ adds nothing there. With $q = 1 - p_{i} > 0$ and $s = k\tau^{q}/q$, the gradient term against the friction scale is $k^{2}\tau^{-2p_{i}}/\tau^{-2} = k^{2}\tau^{2q} = q^{2}s^{2}$, which tends to $0$ as $\tau \to 0$. The exponents are rational on the BKL circle $p(u) = (-u, 1 + u, u(1 + u))/(1 + u + u^{2})$ for rational $u$, the parametrisation of Theorem \claimref{synthesis.bkl-bounce-invariant-holonomy-and-bound}, whose epoch maps $u \mapsto u - 1$ and $u \mapsto 1/(u - 1)$ keep $\sum p = \sum p^{2} = 1$; every exponent on the circle is below $1$ except at the Milne points. The isotropic five-dimensional vacuum point $(\tfrac12, \tfrac12, \tfrac12, -\tfrac12)$ (three observed directions and a contracting hidden one, the direction of Proposition \claimref{synthesis.unit-ball-clauses-hold-at-one-epoch-unless-hidden-direction-co-expands} held static or not) the scalar-shortened isotropic point $(\tfrac13, \tfrac13, \tfrac13)$, with $q_{\phi}^{2} = \tfrac23$, and the interior points $(\tfrac35, \tfrac15, \tfrac15)$ and $(\tfrac12, \tfrac12, 0)$ of the disk have $\sum p = 1$ as well; the points of the disk are those stiff matter reaches, and among the sources of Theorem \claimref{synthesis.kasner-wall-matter-maxwell-bounces-stiff-moves-residue} an aligned Maxwell field is subleading or reflects the weights on the circle by the BKL law and a fluid with $w < 1$ is subleading (its clauses (ii), (iii)), so every epoch those sources reach has $\sum p = 1$; one BKL bounce sequence $u = \tfrac73 \to \tfrac43 \to 3$ keeps $\sum p = 1$ with every exponent below $1$. The check verifies these identities exactly.

\emph{(ii)} With $d/d\tau = k\tau^{-p_{i}}\,d/ds$,
\[
\ddot\phi + \frac{\dot\phi}{\tau} = k^{2}\tau^{-2p_{i}}\phi_{ss} - p_{i}k\tau^{-p_{i}-1}\phi_{s} + k\tau^{-p_{i}-1}\phi_{s} = k^{2}\tau^{-2p_{i}}\Bigl(\phi_{ss} + \frac{q}{k\tau^{q}}\phi_{s}\Bigr) = k^{2}\tau^{-2p_{i}}\Bigl(\phi_{ss} + \frac{1}{s}\phi_{s}\Bigr),
\]
so the mode equation is $k^{2}\tau^{-2p_{i}}(\phi_{ss} + \phi_{s}/s + \phi) = 0$, Bessel's equation of order $0$. The check verifies the operator identity exactly on the terms $s^{n}(\ln s)^{m}$, $n \in \{0, 1, 2, 3, \tfrac52\}$, $m \le 2$, for every exponent of every listed set, with powers of $\tau$ carried as rational exponents; with an added order $\nu^{2} = \tfrac14$ the identity fails (control). In the standard form $v = z\phi$, $v'' + (k^{2} - z''/z)v = 0$ in the axis' conformal time $\eta_{i} = \int d\tau/a_{i} = \tau^{q}/q$, the pump is $z^{2} = \sqrt{|g|}/a_{i} = \tau^{1 - p_{i}} = q\,\eta_{i}$, so $z \propto \eta_{i}^{1/2}$, $z''/z = -\tfrac14\eta_{i}^{-2}$ and the order is $|\tfrac12 - \tfrac12| = 0$. Only $\sum_{j}p_{j} = 1$ was used. (A control in the check: in a de Sitter stage with a hidden direction $b \propto a^{\beta}$ the order is $\tfrac32 + \tfrac{\beta}{2}$ and the index $-\beta$; such a stage can tilt but, as stored, the frozen phase carries no observed mode out.)

\emph{(iii)} The canonical normalisation is $\tau(\phi\dot\phi^{*} - \phi^{*}\dot\phi) = I_{4}$ per mode (the momentum conjugate to $\phi$ is $\sqrt{|g|}\dot\phi = \tau\dot\phi$). For $\phi = NH_{0}^{(1)}(s)$, with $ds/d\tau = qs/\tau$ and $W[H_{0}^{(1)}, H_{0}^{(2)}](s) = -2I_{4}\,W[J_{0}, Y_{0}](s) = -4I_{4}/(\pi s)$, this is $qs|N|^{2}\cdot 4/(\pi s) = 1$ in modulus, so $|N|^{2} = \pi/(4q)$, free of $k$ and $\tau$. The check verifies $W[J_{0}, Y_{0}] = 2/(\pi s)$ exactly on power series to order $s^{21}$, from $J_{0} = \sum_{k}(-1)^{k}(s/2)^{2k}/(k!)^{2}$ and $Y_{0} = \tfrac{2}{\pi}\bigl[(\ln\tfrac{s}{2} + \gamma)J_{0} + \sum_{k \ge 1}(-1)^{k+1}H_{k}(s/2)^{2k}/(k!)^{2}\bigr]$ (standard, cited from memory; the constant multiple of $J_{0}$ does not enter a Wronskian), with a control. Hence $\Delta^{2}(k) = k^{3}|\phi_{k}|^{2}/(2\pi^{2}) = k^{3}M(s)^{2}/(8\pi q)$, $M^{2} = J_{0}^{2} + Y_{0}^{2}$, and at fixed $\tau$, $d\ln s/d\ln k = 1$, so $n - 1 = d\ln\Delta^{2}/d\ln k = 3 + s\,d\ln M^{2}/ds$. By Nicholson's integral $M_{\nu}^{2}(x) = \tfrac{8}{\pi^{2}}\int_{0}^{\infty}K_{0}(2x\sinh t)\cosh 2\nu t\,dt$, $M_{0}^{2}$ is decreasing and, for $\nu < \tfrac12$, $xM_{\nu}^{2}(x)$ is increasing (Watson, cited from memory), so $-1 \le s\,d\ln M^{2}/ds \le 0$ and $n - 1 \in [2, 3]$, with the limits $2$ ($M^{2} \sim 2/(\pi s)$ inside) and $3$. For small $s$, $M^{2} = \tfrac{4}{\pi^{2}}(L^{2} + \tfrac{\pi^{2}}{4}) + O(s^{2}\ln^{2}s)$, $L = \ln\tfrac{s}{2} + \gamma$, so $n - 1 = 3 + 2L/(L^{2} + \pi^{2}/4) \ge 3 + 2/L$; for $s \le 10^{-2}$, $L \le -\ln 200 + \gamma < -5.29 + 0.578 < -4.7$ (with $\ln 200 = 7\ln 2 + \ln 1.5625 > 7(0.693) + 0.446$ and $\gamma < 0.578$), so $n - 1 > 3 - 2/4.7 > 2.57$; at $s$ near $10^{-10}$, $L \approx -23.6$ and $n - 1 > 2.91$. The check samples the leading form on $L \le -4.7$ at both brackets $\pi \in (3, \tfrac{22}{7})$.

\emph{(iv)} With $n_{s} - 1 = -0.0351 \pm 0.0042$ (Fact \claimref{synthesis.imported-planck-2018-lcdm-parameters}), the index is blue by more than $2 - (-0.0351) = 2.0351$. The finite-range bound of Theorem \claimref{synthesis.finite-range-descent-deposit-cannot-seed-cmb}, $\Delta^{2}(k) \le \tfrac{2}{3\pi}\langle\zeta^{2}\rangle(kR)^{3}$, has the same leading power $k^{3}$, but the Kasner spectrum carries the factor $(\ln\tfrac{s}{2} + \gamma)^{2} + \tfrac{\pi^{2}}{4}$, which grows without bound as $k \to 0$ at fixed $\tau$, so the ratio of the spectrum to $(kR)^{3}$ is unbounded for every finite $R$: the exited field is log-correlated, not a deposit of finite range; only the leading power is shared. Time-reversed, a stage expanding away from a wall has $s$ growing, so modes enter and none exits. \emph{Across bounces (argued):} once a mode is outside, the gradient term is negligible, $\tau\dot\phi$ is conserved to leading order and the field shifts by a $k$-independent multiple of $\ln\tau$ in each later epoch, so the $k$-dependence, and with it the index, is carried across a BKL bounce; this is not proved here. \emph{Scope:} a massless minimally coupled spectator scalar (any $\xi R\phi^{2}$ on vacuum points). A field with a dilaton or modulus coupling $f(\varphi)(\partial\chi)^{2}$, among them the Kaluza-Klein vector and the moduli of a hidden direction, has pump $z^{2} = f\,\sqrt{|g|}/a_{i}$, no longer linear in $\eta_{i}$, and its order is not fixed by $\sum p = 1$ (pre-big-bang axions are the standard example of flat or red spectra this way, cited from memory); such fields, and the conversion of a spectator spectrum into the curvature perturbation, are not covered. Momenta along principal axes only (for a general direction the gradient term is $\sum_{j}k_{j}^{2}\tau^{-2p_{j}}$ and the axis of largest $p_{j}$ dominates at the wall, argued); massive fields, whose mass term is negligible against $\tau^{-2}$ at the wall, are argued to give the same order.
